\section{Personlige Refleksioner}
\subsection{Jonathan}
% jonathan
Jeg tog en mindre rolle i projektets opstart og dette påvirkede mit engagement i de tidlige faser af projektet, og jeg oplevede også perioder, hvor jeg ikke deltog i møder eller informerede gruppen om fravær på forhånd. Set i bakspejlet kunne jeg have bidraget mere aktivt. Jeg tror, at min manglende deltagelse i starten kan have været en form for social loafing, hvor jeg ubevidst trak mig tilbage fra ansvaret, fordi jeg følte, at andre i gruppen ville tage sig af opgaverne, eller allerede havde taget sig af dem. For at bryde ud af denne dynamik forsøgte jeg at bidrage på områder, hvor jeg oplevede, at der manglede arbejdskraft, blandt andet gennem kildearbejde og rapportskrivning. Da jeg som person har en tendens til at være tilbageholdene i sociale sammenhænge, var dette en måde jeg stadig kunne bidrage på.
Jeg oplevede at gruppen håndterede situationen professionelt og inkluderende, hvilket gjorde det lettere for mig at bidrage mere aktivt senere i projektforløbet. Jeg har lært, at det er vigtigt at kommunikere åbent om ens deltagelse og forpligtelser i et gruppeprojekt for at sikre, at alle medlemmer føler sig ansvarlige og engagerede. I fremtidige projekter vil jeg stræbe efter at være mere proaktiv i min deltagelse og kommunikere klart med mine gruppemedlemmer om mine bidrag og eventuelle udfordringer, jeg måtte have.
\newpage

\subsection{Victor}
I forhold til selve implementeringsfasen af projektet endte jeg med at tage en meget specifik vinkel på projektet, nemlig implementeringen af voicechat-systemet. Dette endte dog med at være det eneste, jeg bidrog med i forhold til implementering, grundet uforudsete komplikationer. Set i bakspejlet burde jeg have været bedre til at bede om hjælp til visse aspekter af arbejdet, således at jeg kunne have arbejdet på flere dele af spillet frem for kun én enkelt.

Jeg har til tider været sent ude med mine svar i gruppen og missede også et enkelt scrum-møde, hvilket er beklageligt. Især mod slutningen af projektet låste jeg mig fast på at få voicechat-systemet til at fungere, hvilket førte til forsinkede GitHub-pushes, da jeg havde et stort fokus på, at det system, jeg arbejdede på, skulle være fuldt funktionelt.

Den vigtigste læring, jeg tager med fra dette projekt, er især ikke at hyperfiksere på et bestemt problem, men i stedet at blive bedre til at sprøge om hjælp til- og fordele større arbejdsopgaver blandt gruppemedlemmer. På den måde kan jeg fremover få større indflydelse på flere dele af udviklingen og dermed også opnå et bedre indblik i projektets forskellige systemer.

\newpage

\subsection{Nikolaj}
